\chapter*{Resumen}

La electroencefalografía (EEG) constituye una herramienta fundamental en neurociencia clínica y experimental, pero sufre de la pérdida frecuente de canales debido a fallas de contacto, artefactos de movimiento o saturación de amplificadores. La interpolación de estos canales faltantes es un paso crítico del preprocesamiento, y tradicionalmente se ha abordado mediante métodos geométricos como los \textit{spherical splines}, implementados en bibliotecas de referencia como MNE-Python. Sin embargo, estos métodos operan exclusivamente en el dominio espacial, ignorando la conectividad funcional real del cerebro y la dinámica temporal de las señales.

Esta tesis presenta una evaluación comparativa reproducible y exhaustiva de dieciocho métodos de interpolación de canales EEG —geométricos, estadísticos y basados en Procesamiento de Señales en Grafos (GSP)— bajo un protocolo experimental riguroso. El diseño incluye tres bases de datos heterogéneas (PhysioNet de 64 canales, BCI Competition IV 2a de 22 canales y MNE Sample de 60 canales), ciento veinte escenarios de pérdida que combinan modos aleatorio y contiguo con severidades del 10\% al 40\%, y una batería de seis métricas que cubren error de amplitud (MAE, RMSE, SNR), fidelidad morfológica (DTW) y preservación espectral (LSD, Coherencia). Para garantizar una comparación justa, todos los métodos —incluyendo las referencias clásicas— fueron optimizados mediante Optuna con más de quince mil configuraciones paramétricas exploradas, y se incorporó una validación confirmatoria contra la implementación industrial de MNE-Python.

Los resultados revelan que la regularización temporal marca un punto de inflexión decisivo: los métodos de la familia TRSS (\textit{Temporal Regularized Sobolev Smoothness}) superan consistentemente a las referencias geométricas y a los métodos GSP puramente espaciales, con ganancias de 5 a 10~dB en SNR bajo pérdida severa contigua del 40\%. El análisis del mecanismo subyacente muestra que esta ventaja no proviene de una mejor representación espacial, sino de la restricción de continuidad temporal que actúa como información previa de naturaleza fisiológica. No obstante, en regímenes favorables —montajes densos, pérdida aleatoria igual o inferior al 20\%—, MNE-Python iguala o supera a TRSS en preservación espectral, siendo además entre 50 y 100 veces más rápido. Se concluye que TRSS es el método preferido cuando la preservación morfológica bajo pérdida severa es prioritaria, mientras que MNE sigue siendo la opción óptima para entornos con pérdida moderada y restricciones de tiempo real.

La totalidad del código, las configuraciones experimentales, los registros de optimización y los artefactos de resultados están disponibles en un repositorio abierto que garantiza la reproducibilidad completa de los hallazgos.

\vspace{1em}
\noindent\textbf{Palabras clave:} EEG, canales faltantes, interpolación, Procesamiento de Señales en Grafos (GSP), regularización temporal, TRSS, spherical splines, MNE-Python, optimización bayesiana, Optuna, evaluación comparativa reproducible.